\chapter{Métodos Recentes}

\section{DiD linear com adoção escalonada}

\texttt{lpdid} (DiD por Projeções Locais) e \texttt{eventstudy} estimam efeitos de
tratamento em desenhos com adoção escalonada. \texttt{eventstudy} constrói
dummies de tempo do evento em relação ao tratamento e omite o período de
referência.

\begin{lstlisting}[caption={Event study}]
let m_e = eventstudy(y ~ x, df, time=ano, ref=-1,
                     min=-5, max=5)
print(m_e)
\end{lstlisting}

Os coeficientes de cada dummy de tempo do evento mostram o caminho
médio do resultado em relação ao período de tratamento.

\section{Synthetic DiD e CUPED}

Synthetic DiD combina a ideia de controle sintético com diferenças em
diferenças. \texttt{synthdid} estima pesos de unidade e tempo a partir dos
resultados pré-tratamento e reporta o ATT.

\begin{lstlisting}[caption={Synthetic DiD}]
let m_sd = synthdid(df, y, treated, 10, unit="state", period="year")
print(m_sd)
\end{lstlisting}

CUPED ajusta o resultado com covariáveis pré-tratamento para reduzir a
variância em experimentos randomizados.

\begin{lstlisting}[caption={CUPED}]
let m_c = cuped(y ~ x, df, treated="treated")
print(m_c)
\end{lstlisting}

\section{Regressão MIDAS}

Mixed Data Sampling (Ghysels, Santa-Clara e Valkanov, 2004) regrediz uma
variável de baixa frequência sobre defasagens de alta frequência, com
pesos restritos por um polinômio de Almon.

\begin{lstlisting}[caption={MIDAS}]
let m_midas = midas(y ~ x, df, freq=3, lags=12, poly=2)
print(m_midas)
\end{lstlisting}

É útil para nowcasting: usar preditores mensais para prever um alvo
trimestral.
